自旋1/2 粒子起初处于 \(|+z\rangle\) 态。理想 Stern–Gerlach 装置先选择通过 \(+x\) 通道，再选择通过与 z 轴成角 θ（在 xz 平面、方向 \(\boldsymbol n=(\sin\theta,0,\cos\theta)\)）的 \(+\boldsymbol n\) 通道，最后测量 z 分量。只研究被选中的束流。
(1) 求第一次 \(+x\) 选择的概率；(2) 求第二次 \(+\boldsymbol n\) 选择的条件概率；(3) 求整个三次测量得到 \(+z\) 的联合概率；(4) 与不经过中间测量直接测量 z 的概率比较，解释不同测量顺序为何不同。
